\documentclass[conference]{IEEEtran} % tipo de formato que o artigo vai seguir
\usepackage{cite}
\usepackage{graphicx}
\title{Título do artigo}

\author{
	\IEEEauthorblockN{1\textsuperscript{st} Nome 1}
	\IEEEauthorblockA{
		\textit{departamento/diretoria ex.: DIATINF} \\
		\textit{instituição} \\
		Cidade, País \\
		email
	}
	\and
	
	\IEEEauthorblockN{2\textsuperscript{nd} Nome 2}
	\IEEEauthorblockA{
		\textit{departamento/diretoria ex.: DIATINF} \\
		\textit{instituição} \\
		Cidade, País \\
		email
	}
	\and
	
	\IEEEauthorblockN{3\textsuperscript{rd} Nome 3}
	\IEEEauthorblockA{
		\textit{departamento/diretoria ex.: DIATINF} \\
		\textit{instituição} \\
		Cidade, País \\
		email
	}
}

\begin{document}
	
	\maketitle % imprime o titulo
	
	\begin{abstract}
		Resumo do artigo
	\end{abstract}
	
	\begin{IEEEkeywords}
		palavras-chave
	\end{IEEEkeywords}
	
	\section{Introdução}
	% escrever introdução aqui
	
	\section{Seção2}
	% estrutura para usar figura/imagem (está comentado para não dar erros)
	%\begin{figure}[htbp]
	%\centering
	%\includegraphics[width=0.48\textwidth]{imagem.png}
	%\caption{Descrição da figura}
	%\label{fig:modelo}
	%\end{figure}
	% usa \ref{fig:modelo} para citar a figura/img
	% \cite{ref}
	
	\begin{thebibliography}{00}
		\bibitem{ref1} Autor, A. \textit{Título}. Conferência, Ano.
	\end{thebibliography}
\end{document}